%HOMEWORK template for M 325K
\documentclass{article}
\headheight=7pt         \topmargin=14pt
\textheight=574pt       \textwidth=445pt
\oddsidemargin=18pt     \evensidemargin=18pt

%Begin package import

\usepackage{amsmath, amssymb, amsthm}
\usepackage{mathtools}
\usepackage{pdfpages}
\usepackage{tikz-cd}

%End package import
%Begin custom commands

\newcommand{\C}{\mathbb{C}}
\newcommand{\h}[1]{\hat{#1}}
\newcommand{\R}{\mathbb{R}}
\newcommand{\Q}{\mathbb{Q}}
\newcommand{\Z}{\mathbb{Z}}
\newcommand{\N}{\mathbb{N}}
\newcommand{\F}{\mathbb{F}}
\newcommand{\abs}[1]{\left\lvert #1\right\rvert}
\newcommand{\aabs}[1]{\left\lvert \left\lvert #1 \right\rvert\right\rvert}

\newcommand{\RM}[1]{\mathrm{#1}}
\newcommand{\BF}[1]{\textbf{#1}}
\newcommand{\e}{\epsilon}
\newcommand{\ceil}[1]{\lceil{#1}\rceil}
\newcommand{\floor}[1]{\lfloor{#1}\rfloor}
\newcommand{\rarr}[0]{\rightarrow}
\newcommand{\IM}[1]{\text{Im}(#1)}
\newcommand{\Hom}[0]{\text{Hom}}
\newcommand{\Pow}[1]{\text{Pow}(#1)}
\newcommand{\angles}[1]{\langle #1 \rangle}

%End custom commands
%Begin custom theorems

\newtheorem{innercustomgeneric}{\customgenericname}
\providecommand{\customgenericname}{}
\newcommand{\newcustomtheorem}[2]{%
\newenvironment{#1}[1]
{%
\renewcommand\customgenericname{#2}%
\renewcommand\theinnercustomgeneric{##1}%
\innercustomgeneric%
}
{\endinnercustomgeneric}
}

\newcustomtheorem{THRM}{Theorem}
\newcustomtheorem{LEM}{Lemma}
\newcustomtheorem{EX}{Exercise}
\newcustomtheorem{COR}{Corollary}
\newcustomtheorem{PROB}{Problem}
\newcustomtheorem{claim}{Claim}

\newenvironment{solution}
{\renewcommand\qedsymbol{$\blacksquare$}\begin{proof}[Solution]}
{\end{proof}}

\newenvironment{definition}
{\renewcommand\qedsymbol{$\blacksquare$}\begin{proof}[Definition]}
{\end{proof}}

%End custom theorems
\begin{document}

\title{Chapter 2}
\author{Arham Lodha}%put your name here
\date{June 06, 2024}%enter the due date here
\maketitle

\section*{2.6}
\begin{claim}{1}\label{Claim 1.}
    For any summable function f, the Fourier transform 

    \begin{equation*}
        \h{f}(x) = \int f(x) e^{-2 \pi i n x} dx
    \end{equation*}
    exists as an ordinary Lebesgue integral and enjoys the following properties

    \begin{enumerate}
        \item $\aabs{\h{f}}_{L^{\infty}} \leq \aabs{f}_{L^{1}}$
        \item $f \in \mathbf{C}(\R^1)$
        \item $\lim_{\abs{n} \to \infty} \h{f}(n) = 0$  
    \end{enumerate}
    
\end{claim}
\begin{proof}
    \begin{enumerate}
        \item \begin{align*}
            \aabs{\h{f}}_{L^{\infty}} = \max_{n \in \R}\abs{\h{f}(n)} &= \max_{n \in \R} \abs{\int_{-\infty}^{\infty} f(x) e^{-2 \pi i n x} dx} \\
            &\leq \max_{n \in \R} \int_{-\infty}^{\infty} \abs{f(x) e^{-2 \pi i n x}} dx \\
            &= \int_{-\infty}^{\infty} \abs{f(x)} dx \\
            &= \aabs{f}_{L^{1}}
        \end{align*}
        \item \begin{align*}
            \abs{\h{f}(m) - \h{f}(n)} &= \abs{\int_{-\infty}^{\infty}f(x)(e^{-2 \pi i m x} - e^{-2 \pi i n x}) dx} \\
            &\leq \int \abs{f(x)} \abs{e^{-2 \pi i m x} - e^{-2 \pi i n x}} dx \to 0
        \end{align*}

        as $\abs{m - n} \to 0$ because $\lim_{m \to n} \abs{f(x)} \abs{e^{-2 \pi i m x} - e^{-2 \pi i n x}} = 0$ and $\abs{e^{-2 \pi i m x} - e^{-2 \pi i n x}} < \sqrt{2}$ and since $f(x)$ is summable $\abs{f(x)} < M$ thus $\abs{f(x)} \abs{e^{-2 \pi i m x} - e^{-2 \pi i n x}} < M \sqrt{-2}$ for any choice of $m$ and $n$. Thus by dominated convergence $\abs{\h{f}(m) - \h{f}(n)} \to 0$. Thus $\h{f} \in \mathbf{C}(\R)$        

        \item
        Since
        
        \begin{equation*}
            \h{f}(\xi) = -\int f(x) e^{-2 \pi i \xi(x - y)} dx = -\int f(x + y) e^{-2 \pi i \xi x} dx
        \end{equation*}

        for $y = \frac{1}{2\xi}$.
        
        \begin{align*}
            \h{f}(\xi) &= \frac{1}{2}\left[\int f(x) e^{-2 \pi i \xi x} dx - \int f(x + y) e^{-2 \pi i \xi x} dx\right] \\
            &= \frac{1}{2}\int e^{-2 \pi i \xi x}(f(x) - f(x + y)) dx
        \end{align*}

        Thus,

        \begin{equation*}
            \abs{\h{f}(\xi)} = \abs{\frac{1}{2}\int e^{-2 \pi i \xi x}(f(x) - f(x + y)) dx} \leq \frac{1}{2} \int \abs{ e^{-2 \pi i \xi x}(f(x) - f(x + y))}dx = \frac{1}{2}\aabs{f - f_{y}}_{L^1}
        \end{equation*}
        
        As $\abs{\xi} \to \infty \implies y \to 0 \implies \frac{1}{2}\aabs{f - f_{y}}_{L^1} \to 0$ by continuity of translation. Thus $\abs{\h{f}(\xi)} \to 0$ as $\abs{\xi} \to \infty$.    
    \end{enumerate}
\end{proof}

\end{document}